\documentclass{beamer}

\usepackage[utf8]{inputenc}
\usepackage[T1]{fontenc}
\usepackage[brazil]{babel}
\usepackage{graphicx}
\usepackage{booktabs}
\usepackage{listings}
\usepackage{hyperref}
\usepackage{tikz}

\usetheme{Madrid}
\usecolortheme{beaver}
\setbeamertemplate{navigation symbols}{}
\setbeamertemplate{footline}[frame number]

\lstset{
    language=C++,
    basicstyle=\ttfamily\scriptsize,
    keywordstyle=\color{blue}\bfseries,
    breaklines=true
}

\title{Problema da Árvore Geradora com Número Mínimo de Rótulos}
\subtitle{Trabalho Prático -- DCC059 | Teoria dos Grafos}

\author{
    Rian César Quintanilha \and
    Isaac Dizolele Kapela João \and
    Vítor Assunção dos Santos \and
    Welder Salvador Victor
}

\institute{Universidade Federal de Juiz de Fora (UFJF)}
\date{Julho de 2026}

\begin{document}

% ============================================================
% CAPA
% ============================================================

\begin{frame}
    \titlepage
\end{frame}

% ============================================================
% ROTEIRO
% ============================================================

\begin{frame}{Roteiro}
    \tableofcontents
\end{frame}

% ============================================================
\section{Divisão das Tarefas}
% ============================================================

\begin{frame}{Divisão das Tarefas}

\begin{table}[h]
\centering
\renewcommand{\arraystretch}{1.5}
\begin{tabular}{p{6.8cm} p{3.8cm}}
\toprule
\textbf{Atividade} & \textbf{Responsável} \\
\midrule
Estrutura do grafo (lista de adjacência, CRUD) & Welder e Vítor \\
\addlinespace
Heurística gulosa para o MLSTP & Welder e Vítor \\
\addlinespace
Union-Find, remoção de ciclos e implementação do MLSTP & Rian \\
\addlinespace
Geração das 20 instâncias de teste & Isaac \\
\bottomrule
\end{tabular}
\end{table}

\end{frame}

% ============================================================
\section{Definição do Problema}
% ============================================================

\begin{frame}{Definição do Problema}

\begin{columns}[T]

\column{0.55\textwidth}
    Seja $G = (V, E)$ um grafo conexo não direcionado onde cada aresta $e \in E$ possui um rótulo $\ell(e) \in L$.

    \vspace{0.4cm}

    O \textbf{MLSTP} consiste em encontrar uma árvore geradora $T$ que minimize o número de rótulos distintos:

    \[
        \min \left|\{\ell(e) \mid e \in E_T\}\right|
    \]

    \begin{alertblock}{Complexidade}
        NP-difícil — heurísticas são necessárias para instâncias grandes.
    \end{alertblock}

\column{0.42\textwidth}
    \begin{block}{Aplicações}
        \begin{itemize}
            \item Redes de comunicação com múltiplas tecnologias
            \item Redes de transporte
            \item Distribuição de energia
            \item Infraestrutura com restrições de compatibilidade
        \end{itemize}
    \end{block}

\end{columns}

\end{frame}

% ============================================================
\section{Estrutura de Dados}
% ============================================================

\begin{frame}{Estrutura de Dados: Grafo}

\begin{columns}[T]

\column{0.48\textwidth}
    \begin{block}{Classe \texttt{No} (vértice)}
        \begin{itemize}
            \item \texttt{id} — identificador
            \item \texttt{visitado} — flag para buscas
            \item \texttt{grau\_entrada}, \texttt{grau\_saida}
            \item \texttt{vector<No*> vizinhos}
        \end{itemize}
    \end{block}

    \vspace{0.3cm}

    \begin{block}{Struct \texttt{Edge}}
        \begin{itemize}
            \item \texttt{int u, v} — extremos
            \item \texttt{double label} — rótulo
        \end{itemize}
    \end{block}

\column{0.48\textwidth}
    \begin{block}{Classe \texttt{Graph}}
        \begin{itemize}
            \item \texttt{vector<No*> nos}
            \item \texttt{bool direcionado}
        \end{itemize}
    \end{block}

    \vspace{0.3cm}

    \begin{exampleblock}{Representação}
        Lista de adjacência — os rótulos são armazenados na \texttt{struct Edge} e associados às arestas durante a construção do grafo.
    \end{exampleblock}

\end{columns}

\end{frame}

% ============================================================

\begin{frame}{Métodos do Grafo}

\begin{columns}[T]

\column{0.48\textwidth}
    \begin{block}{Operações}
        \begin{itemize}
            \item \texttt{insertVertex / removeVertex}
            \item \texttt{insertEdge / removeEdge}
            \item \texttt{existEdge / changeWeight}
        \end{itemize}
    \end{block}

    \vspace{0.2cm}

    \begin{block}{Consultas}
        \begin{itemize}
            \item \texttt{verifyAdjacency}
            \item \texttt{calculateDegree}
            \item \texttt{listNeighbors}
        \end{itemize}
    \end{block}

\column{0.48\textwidth}
    \begin{block}{Algoritmos}
        \begin{itemize}
            \item \texttt{contructGraphFromEdges()} — constrói o grafo a partir de arestas rotuladas
            \item \texttt{cutCycles()} — remove ciclos via Union-Find
            \item \texttt{minimumLabelingSpanningTree()} — heurística principal do MLSTP
        \end{itemize}
    \end{block}

\end{columns}

\end{frame}

% ============================================================
\section{Union-Find}
% ============================================================

\begin{frame}{Union-Find (Disjoint Set Union)}

\begin{columns}[T]

\column{0.48\textwidth}
    \begin{block}{\texttt{find(x)}}
        Retorna a raiz do conjunto de $x$.\\[0.2cm]
        \textbf{Otimização:} path compression — aponta o nó direto para a raiz, achatando a árvore em chamadas futuras.
    \end{block}

    \vspace{0.3cm}

    \begin{block}{\texttt{unite(x, y)}}
        Une os conjuntos de $x$ e $y$.\\[0.2cm]
        Retorna \texttt{false} se já pertencem ao mesmo conjunto $\Rightarrow$ \textbf{ciclo detectado}.\\[0.2cm]
        \textbf{Otimização:} union by rank — une a árvore menor sob a maior.
    \end{block}

\column{0.48\textwidth}
    \begin{exampleblock}{Uso no projeto}
        \textbf{\texttt{cutCycles()}}\\
        Percorre as arestas do subgrafo gerado pela heurística e remove as que formariam ciclo, garantindo que o resultado final seja uma árvore geradora válida.
    \end{exampleblock}

    \vspace{0.3cm}

    \begin{alertblock}{Complexidade}
        Quase linear: $O(m \cdot \alpha(n))$, onde $\alpha$ é a função inversa de Ackermann.
    \end{alertblock}

\end{columns}

\end{frame}

% ============================================================
\section{Heurística para o MLSTP}
% ============================================================

\begin{frame}{Heurística Gulosa para o MLSTP}

\begin{columns}[T]

\column{0.48\textwidth}
    \begin{block}{Ideia central}
        A cada passo, escolher o rótulo que cobre o \textbf{maior número de vértices ainda não cobertos}.
    \end{block}

    \vspace{0.3cm}

    \textbf{Passos}
    \begin{enumerate}
        \item Agrupar arestas por rótulo.
        \item Enquanto houver vértices não cobertos:
        \begin{itemize}
            \item Calcular o ganho de cada rótulo.
            \item Selecionar o de maior ganho.
            \item Inserir suas arestas no subgrafo.
        \end{itemize}
        \item Chamar \texttt{cutCycles()} para eliminar ciclos.
    \end{enumerate}

\column{0.48\textwidth}
    \begin{exampleblock}{Resultado}
        Subgrafo $\tilde{G} = (V,\, \tilde{E},\, \tilde{L})$ onde:
        \begin{itemize}
            \item $V$ — todos os vértices
            \item $\tilde{L}$ — rótulos selecionados
            \item $\tilde{E}$ — arestas que formam a árvore geradora
        \end{itemize}
    \end{exampleblock}

    \vspace{0.3cm}

    \begin{alertblock}{Característica}
        Heurística gulosa — não garante ótimo global, mas produz soluções de boa qualidade com baixo custo computacional.
    \end{alertblock}

\end{columns}

\end{frame}

% ============================================================
\section{Geração das Instâncias}
% ============================================================


\begin{frame}{Procedimento de Geração (Xiong, 2005)}

\begin{columns}[T]

\column{0.52\textwidth}

    \textbf{Parâmetros de entrada}
    \begin{itemize}
        \item $n$ — número de vértices
        \item $\ell$ — número de rótulos disponíveis
        \item \textit{seed} — garante reprodutibilidade
    \end{itemize}

    \vspace{0.3cm}

    \textbf{Passos}
    \begin{enumerate}
        \item Calcular $m$ (número de arestas da instância)
        \item Gerar \textbf{árvore geradora aleatória} com $n$ vértices \\ $\rightarrow$ conecta todos com $n-1$ arestas
        \item Sortear arestas adicionais sem repetição até completar $m$
        \item Atribuir rótulo aleatório em $[1, \ell]$ a cada aresta
    \end{enumerate}

\column{0.45\textwidth}

    \begin{alertblock}{Por que a árvore geradora primeiro?}
        Sorteio puramente aleatório de $m$ pares pode deixar vértices isolados.

        \vspace{0.2cm}

        A árvore garante conexidade com $n-1$ arestas. As demais arestas adicionam ciclos, mas nunca desconectam o grafo.
    \end{alertblock}

    \vspace{0.3cm}

    \begin{exampleblock}{Reprodutibilidade}
        Usando \texttt{std::mt19937} com seed fixa, a mesma seed sempre gera a mesma instância.
    \end{exampleblock}

\end{columns}

\end{frame}

% ============================================================

\begin{frame}{As 20 Instâncias Geradas}

\begin{columns}[T]

\column{0.55\textwidth}

\begin{table}[h]
\centering
\renewcommand{\arraystretch}{1.15}
\footnotesize
\begin{tabular}{lccr}
\toprule
\textbf{Instância} & $n$ & $\ell$ & $m$ \\
\midrule
01 -- 02 & 100  & 25  & 990     \\
03 -- 04 & 100  & 25  & 2.475   \\
05       & 100  & 25  & 3.960   \\
\midrule
06 -- 07 & 200  & 50  & 3.980   \\
08 -- 09 & 200  & 50  & 9.950   \\
10       & 200  & 50  & 15.920  \\
\midrule
11 -- 12 & 500  & 125 & 24.950  \\
13 -- 14 & 500  & 125 & 62.375  \\
15       & 500  & 125 & 99.800  \\
\midrule
16 -- 17 & 1000 & 250 & 99.900  \\
18 -- 19 & 1000 & 250 & 249.750 \\
20       & 1000 & 250 & 399.600 \\
\bottomrule
\end{tabular}
\end{table}

\column{0.42\textwidth}

    \begin{block}{Critério de variação}
        \begin{itemize}
            \item 4 tamanhos de grafo: \\ $n \in \{100, 200, 500, 1000\}$
            \item 4 tamanhos de rótulos: \\ $\ell \in \{25, 50, 125, 250\}$
            \item 2 seeds por combinação
        \end{itemize}
    \end{block}

    \vspace{0.2cm}

    \begin{exampleblock}{Escala}
        De 990 arestas (instância pequena, esparsa) até 399.600 arestas (instância grande, densa).
    \end{exampleblock}

\end{columns}

\end{frame}

% ============================================================
\section{Conclusão}
% ============================================================

\begin{frame}{Conclusão}

    \textbf{O que foi entregue}
    \begin{itemize}
        \item Estrutura de grafo completa em lista de adjacência
        \item Union-Find com path compression e union by rank
        \item Heurística gulosa para o MLSTP
        \item Gerador de instâncias seguindo Xiong (2005)
        \item 20 instâncias de teste cobrindo diferentes tamanhos e densidades
    \end{itemize}

\end{frame}

% ============================================================

\begin{frame}
    \begin{center}
        {\LARGE Obrigado!}\\[0.8cm]
        {\large Dúvidas?}\\[0.8cm]

        \begin{columns}[c]
        \column{0.5\textwidth}
            \centering
            \textbf{Repositório}\\[0.2cm]
            \footnotesize\url{https://github.com/isaac-kapela/trabalho-final-grafos}

        \column{0.5\textwidth}
            \centering
            \textbf{Relatório}\\[0.2cm]
            \footnotesize\url{https://www.overleaf.com/project/6a4adfd6087b1cdc622e98cc}
        \end{columns}
    \end{center}
\end{frame}

\end{document}
